% TOP_PROBLEM: {"id": 30006675, "title": "Global Langlands Functoriality Conjecture", "category_id": 1, "category": {"id": 1, "name": "number_theory", "display_name": "Number Theory", "description": "Properties of integers, prime numbers, Diophantine equations.", "slug": "number-theory", "order_index": 1, "created_at": "2026-07-31T15:26:25.670Z"}, "status": "open"}
% FIRST_POSTED: 2026-09-27
% ATTEMPT_STATUS: unresolved
% !TeX program = lualatex
% Research continuation by GPT 6 Astra Ultra, 27 September 2026.
% Root statement and inherited source list preserved; new work remains unresolved.
\documentclass[11pt]{article}
\usepackage[margin=25mm]{geometry}
\usepackage{fontspec}
\setmainfont{Latin Modern Roman}
\usepackage{amsmath,amssymb,mathtools}
\usepackage{unicode-math}
\setmathfont{Latin Modern Math}
\usepackage{xurl,hyperref}
\hypersetup{colorlinks=true,urlcolor=blue}
\setcounter{secnumdepth}{0}
\setlength{\parindent}{0pt}
\setlength{\parskip}{5pt}
\setlength{\emergencystretch}{3em}
\begin{document}
% problemId: 30006675
% Canonical problem ID: 30006675
\section*{\texorpdfstring{Global Langlands Functoriality Conjecture}{Global Langlands Functoriality Conjecture}}\label{30006675}
\subsection{Definitions and mathematical statement}
For a number field $F$, ${\mathbb{A}}_F$ is its ring of adeles. A quasisplit connected reductive group has a Borel subgroup over $F$; its $L$-group incorporates its complex dual group and the Galois action. At an unramified place $v$, the Satake parameter $c(\pi_v)$ is a semisimple conjugacy class in the local $L$-group. For an admissible homomorphism $\rho:{}^LG'\to{}^LG$ over the Galois group, the required weak transfer is
\[
 \forall\pi'\in\operatorname{Aut}(G'),\quad \exists\pi\in\operatorname{Aut}(G),\ \exists S\text{ finite},\quad
 c(\pi_v)=\rho_v\bigl(c(\pi'_v)\bigr)\quad(v\notin S).
\]
Here $\operatorname{Aut}(G)$ denotes automorphic representations in Arthur's inclusive convention, not automorphisms of $G$. The source convention, including admissibility, is retained; no cuspidality of the transfer or uniqueness is imposed.
\par\smallskip\textit{Formulation and concept references:} \href{https://www.claymath.org/library/cw/arthur/pdf/76.pdf}{[S1]}.

\subsection{Short English statement}
Does every allowed map of Langlands $L$-groups produce an automorphic transfer matching the local arithmetic data almost everywhere?
\subsection{Research attempt}

\paragraph{Scope and source checkpoint.}
The selected assertion is weak global transfer: matching Satake classes
outside a finite set, with the inclusive automorphic convention of \href{https://www.claymath.org/library/cw/arthur/pdf/76.pdf}{[S1]}.
Neither cuspidality nor matching every ramified component is added. Arthur
\href{https://www.claymath.org/library/cw/arthur/pdf/76.pdf}{[S1]}, Section 1, supplies this formulation; \href{https://www.math.toronto.edu/arthur/pdf/arthur-func-and-trace-formula.pdf}{[S2]} discusses the role of
global trace identities. The present argument does not supply those
identities or claim a new general transfer theorem. It proves a split-torus
subcase directly and constructs an exact obstruction to assembling a
transfer merely from individually admissible local parameters.

We use id\`ele-class characters as the concrete automorphic objects for
$\mathrm{GL}_1$. Their conventions are reviewed in \href{https://math.mit.edu/~poonen/786/notes.pdf}{[E1]}, Section 5.8.
The distinction between a character of the id\`ele group and a character
of its quotient by the rational points is central to the attempt.

\paragraph{Attempt A: monomial maps between split tori.}
Fix a number field $F$ and split tori
$T'=\mathbb G_m^r$, $T=\mathbb G_m^s$. Consider the $L$-homomorphism
which is the identity on the Galois factor and on dual tori is
\[
 \rho(z_1,\ldots,z_r)_j=\prod_{i=1}^r z_i^{a_{ji}},
 \qquad A=(a_{ji})\in M_{s\times r}(\mathbb Z).
\]
This fixes a genuine subcase of the catalogue: no additional Galois
cocycle in the $L$-homomorphism is being treated. For automorphic
characters $\chi_i$ of $F^\times\backslash\mathbb A_F^\times$, define
\begin{equation}\label{cont7:torus}
 \psi_j=\prod_{i=1}^r\chi_i^{a_{ji}},\qquad
 \pi=\psi_1\boxtimes\cdots\boxtimes\psi_s.
\end{equation}
Products and inverses preserve continuity and triviality on $F^\times$.
Consequently these are again automorphic characters. In the usual
smooth automorphic convention they have the required local smoothness
and archimedean character behavior as well. If the source characters
are unitary, so are the target characters; unitarity is not needed to
replace the inclusive convention of the question.

Let $S$ contain the finitely many places where one of the $\chi_i$ is
ramified. At a finite $v\notin S$, choose a uniformizer $\varpi_v$ and
use the same reciprocity convention on source and target. Then
\[
 c(\pi_v)_j=\psi_{j,v}(\varpi_v)
       =\prod_i\chi_{i,v}(\varpi_v)^{a_{ji}}
       =\rho\bigl(c(\pi'_v)\bigr)_j.
\]
If the alternative inverse-uniformizer convention is chosen, all the
parameters invert and the same identity holds. Thus the required
almost-everywhere identity follows without making a local-to-global
existence assumption: the global characters were constructed first.

The construction also respects composition. For an integer matrix $B$
describing a second torus map, applying it to \eqref{cont7:torus} gives
exponents $\sum_jb_{kj}a_{ji}$, exactly the matrix $BA$ describing the
composite map. Zero rows give the trivial character; negative entries
give inverse characters. Repeated characters and kernels of the torus
map cause no problem. This checks several degenerate cases that a
claim requiring cuspidal output would mishandle.

\paragraph{Where the same construction stops for general groups.}
For a general $\rho$, one can still write the desired semisimple
conjugacy classes $\rho_v(c(\pi'_v))$. One may also form their local
Euler factors, after choosing a representation of the dual target
group. Neither operation constructs a vector in an automorphic
representation of $G$. In particular, the fact that each local
parameter belongs to the correct dual group imposes no visibly global
summation identity. The following test makes this distinction exact
already in dimension one.

\paragraph{Lemma 7.1: a rigidity test for rational Hecke characters.}
Suppose a continuous id\`ele-class character of $\mathbb Q$ has
unramified parameters in $\{1,-1\}$ at all but finitely many primes.
Then those parameters agree almost everywhere with a finite-order
Dirichlet character.

\emph{Proof.} The elementary id\`ele-class description is
\[
 \mathbb Q^\times\backslash\mathbb A_{\mathbb Q}^\times
           \simeq\mathbb R_{>0}\times\widehat{\mathbb Z}^{\times}.
\]
To see the representative explicitly, multiply an id\`ele by the
rational number cancelling each of its finitely many nonzero finite
valuations, then use a rational sign to make its real component
positive. The finite components are now units, and no nontrivial
rational multiplication preserves these normalizations. This also
gives the stated product topology.

A continuous complex character of $\mathbb R_{>0}$ is $x\mapsto x^u$
for $u\in\mathbb C$. A continuous character of
$\widehat{\mathbb Z}^{\times}$ has finite image: choose an open subgroup
mapped into a sufficiently small arc about one in the circle; that arc
contains no nontrivial subgroup, so the open subgroup is in the kernel.
Here compactness first ensures that the image has modulus one.
The finite character therefore factors through
$(\mathbb Z/m\mathbb Z)^\times$ for some $m$. Translating to prime
parameters, and absorbing inverses into the conventions, gives
\begin{equation}\label{cont7:hecke}
 \alpha_p=\chi(p)p^u\qquad(p\nmid m)
\end{equation}
for a finite-order Dirichlet character $\chi$.

Since $|\alpha_p|=1$, $\operatorname{Re}u=0$; write $u=it$. If $h$ is
the order of $\chi$, then for every prime outside the exceptional set
\[
 p^{2hit}=1,\qquad ht\log p\in\pi\mathbb Z.
\]
Choose two distinct such primes $p,q$. If $t\ne0$, both corresponding
integers are nonzero, so $\log p/\log q$ is rational. Clearing the
denominator would give an equality between a positive power of $p$
and a positive power of $q$, contrary to unique factorization. Thus
$t=0$, proving the lemma.\hfill$\square$

This proof needs no assertion about primes in arithmetic progressions.
Only two distinct primes outside a finite set are required. It also
allows the putative global character to have arbitrary finite
conductor and an initially unrestricted archimedean component.

\paragraph{Attempt B: diagonalize against all possible global characters.}
Enumerate all Dirichlet characters as $\chi_1,\chi_2,\ldots$, retaining
a modulus $m_j$ for each. Enumerate all pairs $(j,k)$ of positive
integers in a sequence. At the stage assigned to $(j,k)$, choose a
previously unused prime $p_{j,k}$ not dividing $m_j$, and assign
\[
 \alpha_{p_{j,k}}\in\{1,-1\}\setminus\{\chi_j(p_{j,k})\}.
\]
There is always such a sign. If the character value is neither sign,
either sign works. At every finite stage only finitely many primes are
excluded, so a new prime is available. Assign $\alpha_p=1$ at any
primes not selected by this process.

\textbf{Proposition 7.2.} Every $\alpha_p$ is the parameter of an
unramified unitary local character of $\mathbb Q_p^\times$, but the
system $(\alpha_p)_p$ is not the almost-everywhere parameter system of
any automorphic character of $\mathrm{GL}_1/\mathbb Q$.

\emph{Proof.} Locally define
$\theta_p(x)=\alpha_p^{v_p(x)}$. This is continuous, trivial on units,
and unitary. Suppose a global character matched these parameters
outside a finite set. Lemma 7.1 would identify its parameters there
with one of the enumerated $\chi_j$. But the construction makes them
different at every $p_{j,k}$, an infinite set of distinct primes. A
finite exceptional set cannot remove all those disagreements.
\hfill$\square$

One can even form a continuous character of the full id\`ele group by
the product of these local characters and a trivial real component:
the product is finite on each id\`ele because almost all its finite
components are units. Its missing property is descent to the quotient
by $\mathbb Q^\times$. Changing finitely many local factors and
choosing a different archimedean character cannot repair that defect,
as Proposition 7.2 proves.

\paragraph{Stress test: an Euler product is not an automorphy certificate.}
For the same parameters the formal product
\[
 L_\alpha(s)=\prod_p(1-\alpha_pp^{-s})^{-1}
\]
is an actual absolutely convergent, nonvanishing holomorphic function
in $\operatorname{Re}s>1$. Indeed its logarithm is the absolutely
convergent sum
\[
 \sum_p\sum_{a\ge1}\frac{\alpha_p^a}{a p^{as}},
\]
whose absolute value series is bounded on each
$\operatorname{Re}s\ge1+\varepsilon$ by a constant times
$\sum_{n\ge2}n^{-1-\varepsilon}$.
Thus correct degree-one local factors, unit-modulus parameters,
multiplicativity, and absolute convergence in a right half-plane can
all coexist with failure of global automorphy. No analytic
continuation or functional equation is claimed for this product.
Those would be substantial extra conditions, not consequences of
the displayed convergence calculation.

This example is \emph{not} a counterexample to functoriality. It has
not been obtained by applying an admissible $\rho$ to an automorphic
source. What it disproves is the intermediate assertion that arbitrary
well-formed local parameters can automatically be assembled into the
needed global object. A proof must use the source automorphy in an
essential way beyond computing the transformed Euler factors.

\paragraph{Conditional continuation and first unresolved step.}
For any fixed source and $\rho$, a successful continuation could
construct a nonzero automorphic eigenrepresentation whose spherical
Hecke eigenvalues equal the characters determined by
$\rho_v(c(\pi'_v))$ outside a finite set. If such a representation is
constructed with the catalogue's automorphic convention, the defining
Satake correspondence gives the required transfer. What is missing is
its global existence, not the formal description of its prospective
eigenvalues. Producing compatible finite lists of local eigenvalues
does not prove that one automorphic representation realizes all of
them; no compactness theorem for that assertion has been established
here. In approaches through $L$-functions one would likewise need the
specific analytic and twisting hypotheses of an applicable converse
theorem, with all its scope restrictions proved.

\paragraph{Verification and outcome.}
The torus identity and composition rule are exact character
calculations. The obstruction is an infinite diagonal construction
with every choice justified; it is not extrapolated from a finite
prime sample. The rigidity argument explicitly handles both the
finite exceptional set and the archimedean norm twist, two freedoms
that are essential to the weak-transfer question.

\textbf{Outcome: Proved split-torus subcase and a precise failure of
local assembly; UNRESOLVED.} The general reductive-group transfer
requires a global existence argument using the automorphic source.
Nothing in this attempt supplies that argument or settles general
Langlands functoriality.

\subsection{Sources}
\begin{itemize}
\item[S1]Classifying automorphic representations, Clay collected-works item76.
\url{https://www.claymath.org/library/cw/arthur/pdf/76.pdf}.
\textit{Formulation source.}
\item[S2]Functoriality and the Trace Formula — James Arthur.
\url{https://www.math.toronto.edu/arthur/pdf/arthur-func-and-trace-formula.pdf}.
\textit{Further reference; not a proof endorsement.}
\item[S3]Math 590: Open Problems in Number Theory, Spring 2026 — Duke University.
\url{https://sites.math.duke.edu/~pierce/documents/Math590.pdf}.
\textit{Further reference; not a proof endorsement.}
\item[E1]Bjorn Poonen, \emph{Tate's Thesis}, MIT 18.786 course notes, Section 5.8 (id\`ele-class characters).
\url{https://math.mit.edu/~poonen/786/notes.pdf}.
\textit{Primary author notes consulted 27 September 2026. The rational-character rigidity and diagonal obstruction are proved in this attempt.}
\end{itemize}
\end{document}
